%        File: revise.tex
%     Created: Wed Oct 27 02:00 PM 2018 P
% Last Change: Wed Oct 27 02:00 PM 2018 P
%

%
% Copyright 2007, 2008, 2009 Elsevier Ltd
%
% This file is part of the 'Elsarticle Bundle'.
% ---------------------------------------------
%
% It may be distributed under the conditions of the LaTeX Project Public
% License, either version 1.2 of this license or (at your option) any
% later version.  The latest version of this license is in
%    http://www.latex-project.org/lppl.txt
% and version 1.2 or later is part of all distributions of LaTeX
% version 1999/12/01 or later.
%
% The list of all files belonging to the 'Elsarticle Bundle' is
% given in the file `manifest.txt'.
%

% Template article for Elsevier's document class `elsarticle'
% with numbered style bibliographic references
% SP 2008/03/01
%
%
%
% $Id: elsarticle-template-num.tex 4 2009-10-24 08:22:58Z rishi $
%
%
%\documentclass[preprint,12pt]{elsarticle}
\documentclass[answers,11pt]{exam}

% \documentclass[preprint,review,12pt]{elsarticle}

% Use the options 1p,twocolumn; 3p; 3p,twocolumn; 5p; or 5p,twocolumn
% for a journal layout:
% \documentclass[final,1p,times]{elsarticle}
% \documentclass[final,1p,times,twocolumn]{elsarticle}
% \documentclass[final,3p,times]{elsarticle}
% \documentclass[final,3p,times,twocolumn]{elsarticle}
% \documentclass[final,5p,times]{elsarticle}
% \documentclass[final,5p,times,twocolumn]{elsarticle}

% if you use PostScript figures in your article
% use the graphics package for simple commands
% \usepackage{graphics}
% or use the graphicx package for more complicated commands
\usepackage{graphicx}
% or use the epsfig package if you prefer to use the old commands
% \usepackage{epsfig}

% The amssymb package provides various useful mathematical symbols
\usepackage{amssymb}
% The amsthm package provides extended theorem environments
% \usepackage{amsthm}
\usepackage{amsmath}

% The lineno packages adds line numbers. Start line numbering with
% \begin{linenumbers}, end it with \end{linenumbers}. Or switch it on
% for the whole article with \linenumbers after \end{frontmatter}.
\usepackage{lineno}

% I like to be in control
\usepackage{placeins}

% natbib.sty is loaded by default. However, natbib options can be
% provided with \biboptions{...} command. Following options are
% valid:

%   round  -  round parentheses are used (default)
%   square -  square brackets are used   [option]
%   curly  -  curly braces are used      {option}
%   angle  -  angle brackets are used    <option>
%   semicolon  -  multiple citations separated by semi-colon
%   colon  - same as semicolon, an earlier confusion
%   comma  -  separated by comma
%   numbers-  selects numerical citations
%   super  -  numerical citations as superscripts
%   sort   -  sorts multiple citations according to order in ref. list
%   sort&compress   -  like sort, but also compresses numerical citations
%   compress - compresses without sorting
%
% \biboptions{comma,round}

% \biboptions{}


\usepackage{xspace}
\usepackage{color}

\usepackage{multirow}
\usepackage[hyphens]{url}


\usepackage[acronym,toc]{glossaries}
\newacronym[longplural={metric tons of heavy metal}]{MTHM}{MTHM}{metric ton of heavy metal}
\newacronym{ABM}{ABM}{agent-based modeling}
\newacronym{ACDIS}{ACDIS}{Program in Arms Control \& Domestic and International Security}
\newacronym{AFCI}{AFCI}{Advanced Fuel Cycle Initiative}
\newacronym{AHTR}{AHTR}{Advanced High Temperature Reactor}
\newacronym{ANDRA}{ANDRA}{Agence Nationale pour la gestion des D\'echets RAdioactifs, the French National Agency for Radioactive Waste Management}
\newacronym{ANL}{ANL}{Argonne National Laboratory}
\newacronym{ARDP}{ARDP}{Advanced Reactors Demonstration Program}
\newacronym{API}{API}{application programming interface}
\newacronym{ARCH}{ARCH}{autoregressive conditional heteroskedastic}
\newacronym{ARE}{ARE}{Aircraft Reactor Experiment}
\newacronym{ARFC}{ARFC}{Advanced Reactors and Fuel Cycles}
\newacronym{ARMA}{ARMA}{autoregressive moving average}
\newacronym{ASME}{ASME}{American Society of Mechanical Engineers}
\newacronym{ATWS}{ATWS}{Anticipated Transient Without Scram}
\newacronym{BDBE}{BDBE}{Beyond Design Basis Event}
\newacronym{BIDS}{BIDS}{Berkeley Institute for Data Science}
\newacronym{BOL}{BOL}{Beginning-of-Life}
\newacronym{BSD}{BSD}{Berkeley Software Distribution}
\newacronym{BWR}{BWR}{Boiling Water Reactor}
\newacronym{CAFCA}{CAFCA}{ Code for Advanced Fuel Cycles Assessment }
\newacronym{CASL}{CASL}{Consortium for Advanced Simulation of Light Water Reactors}
\newacronym{CDTN}{CDTN}{Centro de Desenvolvimento da Tecnologia Nuclear}
\newacronym{CEA}{CEA}{Commissariat \`a l'\'Energie Atomique et aux \'Energies Alternatives}
\newacronym{CI}{CI}{continuous integration}
\newacronym{CLASS}{CLASS}{Core Library for Advanced Scenario Simulation}
\newacronym{CNEC}{CNEC}{Consortium for Nonproliferation Enabling Capabilities}
\newacronym{CNEN}{CNEN}{Comiss\~{a}o Nacional de Energia Nuclear}
\newacronym{CNERG}{CNERG}{Computational Nuclear Engineering Research Group}
\newacronym{COSI}{COSI}{Commelini-Sicard}
\newacronym{COTS}{COTS}{commercial, off-the-shelf}
\newacronym{CRAM}{CRAM}{Chebyshev rational approximation method}
\newacronym{CSNF}{CSNF}{commercial spent nuclear fuel}
\newacronym{CTAH}{CTAHs}{Coiled Tube Air Heaters}
\newacronym{CUBIT}{CUBIT}{CUBIT Geometry and Mesh Generation Toolkit}
\newacronym{CURIE}{CURIE}{Centralized Used Fuel Resource for Information Exchange}
\newacronym{CYBORG}{CyBORG}{Cyclus-Based ORiGen}
\newacronym{DAG}{DAG}{directed acyclic graph}
\newacronym{DANESS}{DANESS}{Dynamic Analysis of Nuclear Energy System Strategies}
\newacronym{DBE}{DBE}{Design Basis Event}
\newacronym{DESAE}{DESAE}{Dynamic Analysis of Nuclear Energy Systems Strategies}
\newacronym{DHS}{DHS}{Department of Homeland Security}
\newacronym{DOE}{DOE}{Department of Energy}
\newacronym{DOE-NE}{DOE-NE}{U.S. Department of Energy, Office of Nuclear Energy}
\newacronym{DRACS}{DRACS}{Direct Reactor Auxiliary Cooling System}
\newacronym{DRE}{DRE}{dynamic resource exchange}
\newacronym{DSNF}{DSNF}{DOE spent nuclear fuel}
\newacronym{DYMOND}{DYMOND}{Dynamic Model of Nuclear Development }
\newacronym{EBS}{EBS}{Engineered Barrier System}
\newacronym{EBR}{EBR-II}{Experimental Breeder Reactor II}
\newacronym{EDZ}{EDZ}{Excavation Disturbed Zone}
\newacronym{EG}{EG}{Evaluation Group}
\newacronym{EIA}{EIA}{U.S. Energy Information Administration}
\newacronym{EOL}{EOL}{End-of-life}
\newacronym{EPA}{EPA}{Environmental Protection Agency}
\newacronym{EP}{EP}{Engineering Physics}
\newacronym{ES}{E\&S}{Evaluation and Screening}
\newacronym{FCM}{FCM}{Fully Ceramic Microencapsulated}
\newacronym{FCO}{FCO}{Fuel Cycle Options}
\newacronym{FCT}{FCT}{Fuel Cycle Technology}
\newacronym{FCWMD}{FCWMD}{Fuel Cycle and Waste Management Division}
\newacronym{FEHM}{FEHM}{Finite Element Heat and Mass Transfer}
\newacronym{FEPs}{FEPs}{Features, Events, and Processes}
\newacronym{FHR}{FHR}{Fluoride-Salt-Cooled High-Temperature Reactor}
\newacronym{FLiBe}{FLiBe}{Fluoride-Lithium-Beryllium}
\newacronym{GCAM}{GCAM}{Global Change Assessment Model}
\newacronym{GDSE}{GDSE}{Generic Disposal System Environment}
\newacronym{GDSM}{GDSM}{Generic Disposal System Model}
\newacronym{GEH}{GEH}{GE-Hitachi}
\newacronym{GENIUSv1}{GENIUSv1}{Global Evaluation of Nuclear Infrastructure Utilization Scenarios, Version 1}
\newacronym{GENIUSv2}{GENIUSv2}{Global Evaluation of Nuclear Infrastructure Utilization Scenarios, Version 2}
\newacronym{GENIUS}{GENIUS}{Global Evaluation of Nuclear Infrastructure Utilization Scenarios}
\newacronym{GPAM}{GPAM}{Generic Performance Assessment Model}
\newacronym{GRSAC}{GRSAC}{Graphite Reactor Severe Accident Code}
\newacronym{GUI}{GUI}{graphical user interface}
\newacronym{HALEU}{HALEU}{High Assay Low Enriched Uranium}
\newacronym{HEU}{HEU}{High Enriched Uranium}
\newacronym{HLW}{HLW}{high level waste}
\newacronym{HPC}{HPC}{high-performance computing}
\newacronym{HTC}{HTC}{high-throughput computing}
\newacronym{HTGR}{HTGR}{High Temperature Gas-Cooled Reactor}
\newacronym{HTTR}{HTTR}{High Temperature Engineering Test Reactor}
\newacronym{IAEA}{IAEA}{International Atomic Energy Agency}
\newacronym{IEMA}{IEMA}{Illinois Emergency Mangament Agency}
\newacronym{INL}{INL}{Idaho National Laboratory}
\newacronym{IPRR1}{IRP-R1}{Instituto de Pesquisas Radioativas Reator 1}
\newacronym{IRP}{IRP}{Integrated Research Project}
\newacronym{ISFSI}{ISFSI}{Independent Spent Fuel Storage Installation}
\newacronym{ISRG}{ISRG}{Independent Student Research Group}
\newacronym{JFNK}{JFNK}{Jacobian-Free Newton Krylov}
\newacronym{LACE}{LACE}{Los Alamos Criticality Engine}
\newacronym{LANL}{LANL}{Los Alamos National Laboratory}
\newacronym{LBNL}{LBNL}{Lawrence Berkeley National Laboratory}
\newacronym{LCOE}{LCOE}{levelized cost of electricity}
\newacronym{LDRD}{LDRD}{laboratory directed research and development}
\newacronym{LEU}{LEU}{Low Enriched Uranium}
\newacronym{LFR}{LFR}{Lead-Cooled Fast Reactor}
\newacronym{LGPL}{LGPL}{Lesser GNU Public License}
\newacronym{LLNL}{LLNL}{Lawrence Livermore National Laboratory}
\newacronym{LLW}{LLW}{low level waste}
\newacronym{LMFBR}{LMFBR}{Liquid-Metal-cooled Fast Breeder Reactor}
\newacronym{LOFC}{LOFC}{Loss of Forced Cooling}
\newacronym{LOHS}{LOHS}{Loss of Heat Sink}
\newacronym{LOLA}{LOLA}{Loss of Large Area}
\newacronym{LP}{LP}{linear program}
\newacronym{LWR}{LWR}{Light Water Reactor}
\newacronym{MARKAL}{MARKAL}{MARKet and ALlocation}
\newacronym{MA}{MA}{minor actinide}
\newacronym{MCNP}{MCNP}{Monte Carlo N-Particle code}
\newacronym{MILP}{MILP}{mixed-integer linear program}
\newacronym{MIT}{MIT}{the Massachusetts Institute of Technology}
\newacronym{MMR}{MMR}{Micro Modular Reactor}
\newacronym{MOAB}{MOAB}{Mesh-Oriented datABase}
\newacronym{MOL}{MOL}{Middle-of-Life}
\newacronym{MOOSE}{MOOSE}{Multiphysics Object-Oriented Simulation Environment}
\newacronym{MOX}{MOX}{mixed oxide}
\newacronym{MSBR}{MSBR}{Molten Salt Breeder Reactor}
\newacronym{MSRE}{MSRE}{Molten Salt Reactor Experiment}
\newacronym{MSR}{MSR}{Molten Salt Reactor}
\newacronym{NAGRA}{NAGRA}{National Cooperative for the Disposal of Radioactive Waste}
\newacronym{NCSA}{NCSA}{National Center for Supercomputing Applications}
\newacronym{NEAMS}{NEAMS}{Nuclear Engineering Advanced Modeling and Simulation}
\newacronym{NEUP}{NEUP}{Nuclear Energy University Programs}
\newacronym{NFCSim}{NFCSim}{Nuclear Fuel Cycle Simulator}
\newacronym{NFC}{NFC}{Nuclear Fuel Cycle}
\newacronym{NGNP}{NGNP}{Next Generation Nuclear Plant}
\newacronym{NIST}{NIST}{National Institute of Standards and Technology}
\newacronym{NMWPC}{NMWPC}{Nuclear MW Per Capita}
\newacronym{NNSA}{NNSA}{National Nuclear Security Administration}
\newacronym{NPRE}{NPRE}{Department of Nuclear, Plasma, and Radiological Engineering}
\newacronym{NQA1}{NQA-1}{Nuclear Quality Assurance - 1}
\newacronym{NRC}{NRC}{Nuclear Regulatory Commission}
\newacronym{NSF}{NSF}{National Science Foundation}
\newacronym{NSSC}{NSSC}{Nuclear Science and Security Consortium}
\newacronym{NUWASTE}{NUWASTE}{Nuclear Waste Assessment System for Technical Evaluation}
\newacronym{NWF}{NWF}{Nuclear Waste Fund}
\newacronym{NWTRB}{NWTRB}{Nuclear Waste Technical Review Board}
\newacronym{OAT}{OAT}{one-at-a-time}
\newacronym{OCRWM}{OCRWM}{Office of Civilian Radioactive Waste Management}
\newacronym{OECD}{OECD}{Organisation for Economic Cooperation and Development}
\newacronym{ORNL}{ORNL}{Oak Ridge National Laboratory}
\newacronym{PARCS}{PARCS}{Purdue Advanced Reactor Core Simulator}
\newacronym{PBAHTR}{PB-AHTR}{Pebble Bed Advanced High Temperature Reactor}
\newacronym{PBFHR}{PB-FHR}{Pebble-Bed Fluoride-Salt-Cooled High-Temperature Reactor}
\newacronym{PEI}{PEI}{Peak Environmental Impact}
\newacronym{PH}{PRONGHORN}{PRONGHORN}
\newacronym{PI}{PI}{Principal Investigator}
\newacronym{PNNL}{PNNL}{Pacific Northwest National Laboratory}
\newacronym{PRIS}{PRIS}{Power Reactor Information System}
\newacronym{PRKE}{PRKE}{Point Reactor Kinetics Equations}
\newacronym{PSPG}{PSPG}{Pressure-Stabilizing/Petrov-Galerkin}
\newacronym{PWAR}{PWAR}{Pratt and Whitney Aircraft Reactor}
\newacronym{PWR}{PWR}{Pressurized Water Reactor}
\newacronym{PyNE}{PyNE}{Python toolkit for Nuclear Engineering}
\newacronym{PyRK}{PyRK}{Python for Reactor Kinetics}
\newacronym{QA}{QA}{quality assurance}
\newacronym{RDD}{RD\&D}{Research Development and Demonstration}
\newacronym{RD}{R\&D}{Research and Development}
\newacronym{RELAP}{RELAP}{Reactor Excursion and Leak Analysis Program}
\newacronym{RIA}{RIA}{Reactivity Insertion Accident}
\newacronym{RIF}{RIF}{Region-Institution-Facility}
\newacronym{SAM}{SAM}{Simulation and Modeling}
\newacronym{SCF}{SCF}{Software Carpentry Foundation}
\newacronym{SFR}{SFR}{Sodium-Cooled Fast Reactor}
\newacronym{SINDAG}{SINDA{\textbackslash}G}{Systems Improved Numerical Differencing Analyzer $\backslash$ Gaski}
\newacronym{SKB}{SKB}{Svensk K\"{a}rnbr\"{a}nslehantering AB}
\newacronym{SMR}{SMR}{small modular reactor}
\newacronym{SNF}{SNF}{spent nuclear fuel}
\newacronym{SNL}{SNL}{Sandia National Laboratory}
\newacronym{SNM}{SNM}{Special Nuclear Material}
\newacronym{SRS}{SRS}{Savannah River Site}
\newacronym{STC}{STC}{specific temperature change}
\newacronym{SUPG}{SUPG}{Streamline-Upwind/Petrov-Galerkin}
\newacronym{SWF}{SWF}{Separations and Waste Forms}
\newacronym{SWU}{SWU}{Separative Work Unit}
\newacronym{SandO}{S\&O}{Signatures and Observables}
\newacronym{THW}{THW}{The Hacker Within}
\newacronym{TRIGA}{TRIGA}{Training Research Isotope General Atomic}
\newacronym{TRISO}{TRISO}{Tristructural Isotropic}
\newacronym{TRU}{TRU}{transuranic}
\newacronym{TSM}{TSM}{Total System Model}
\newacronym{TSPA}{TSPA}{Total System Performance Assessment for the Yucca Mountain License Application}
\newacronym{UDB}{UDB}{Unified Database}
\newacronym{UFD}{UFD}{Used Fuel Disposition}
\newacronym{UML}{UML}{Unified Modeling Language}
\newacronym{UNF}{UNF}{used nuclear fuel}
\newacronym{UNFSTANDARDS}{UNFST\&DARDS}{Used Nuclear Fuel Storage, Transportation \& Disposal Analysis Resource and Data System}
\newacronym{USNC}{USNC}{Ultra Safe Nuclear Company}
\newacronym{UOX}{UOX}{uranium oxide}
\newacronym{UQ}{UQ}{uncertainty quantification}
\newacronym{US}{US}{United States}
\newacronym{UW}{UW}{University of Wisconsin}
\newacronym{VISION}{VISION}{the Verifiable Fuel Cycle Simulation Model}
\newacronym{VV}{V\&V}{verification and validation}
\newacronym{WIPP}{WIPP}{Waste Isolation Pilot Plant}
\newacronym{YMG}{YMG}{Young Members Group}
\newacronym{YMR}{YMR}{Yucca Mountain Repository Site}
\newacronym{NEI}{NEI}{Nuclear Energy Institute}
%\newacronym{<++>}{<++>}{<++>}
%\newacronym{<++>}{<++>}{<++>}

\makeglossaries

%\journal{Nuclear Technology}

\begin{document}

%\begin{frontmatter}

% Title, authors and addresses


\title{Effects of uranium impurities in downblended HEU on
HTGR performance}

% Authors

% use the tnoteref command within \title for footnotes;
% use the tnotetext command for the associated footnote;
% use the fnref command within \author or \address for footnotes;
% use the fntext command for the associated footnote;
% use the corref command within \author for corresponding author footnotes;
% use the cortext command for the associated footnote;
% use the ead command for the email address,
% and the form \ead[url] for the home page:
%
% \title{Title\tnoteref{label1}}
% \tnotetext[label1]{}
% \author{Name\corref{cor1}\fnref{label2}}
% \ead{email address}
% \ead[url]{home page}
% \fntext[label2]{}
% \cortext[cor1]{}
% \address{Address\fnref{label3}}
% \fntext[label3]{}

\author{Amanda M. Bachmann, 
Zo\"{e} Richter,
Kathryn D. Huff,
Madicken Munk}

% use optional labels to link authors explicitly to addresses:
% \author[label1,label2]{<author name>}
% \address[label1]{<address>}
% \address[label2]{<address>}


%\author[uiuc]{Kathryn Huff}
%        \ead{kdhuff@illinois.edu}
%  \address[uiuc]{Department of Nuclear, Plasma, and Radiological Engineering,
%        118 Talbot Laboratory, MC 234, Universicy of Illinois at
%        Urbana-Champaign, Urbana, IL 61801}
%
% \end{frontmatter}
\maketitle
\section*{Review General Response}
We would like to thank the reviewers for their detailed assessment of this
paper. Your comments have resulted in changes which certainly improved the
paper.


\section*{Reviewer 1}
\begin{questions}


        %---------------------------------------------------------------------
	\question Biggest observation is that the magnitude of the fluxes for total,
	radial, and axial seem high and inconsistent between the radial and axial
	shapes.  Some of the conclusions about the coolant and moderator temperature
	coefficients could use some more rationale tied to the Serpent output results
	for reaction rates in the regions or flux per unit lethargy plots (with perhaps
	some more methods explanation for the geometric expansion in addition to the
	simple cross section changes for just cross section temperature effects)
	 \begin{solution}
	 \end{solution}
	\end{questions} 


\section*{Reviewer 2}
\paragraph{Overall: } The paper is of significant interest to the community.
However, additional edits are recommended before this paper advances to
publication stage.
The paper represents essentially a very direct approach - (1) simplified
models, (2) a few compositions selected, (3) integral differential metrics for
the analysis, (4) results of the limited analysis space. The paper needs to be
expanded to prove the validity and relevance of the results. As provided,
results appear to be artifacts of the metrics and the modeling approach
limitations.

\begin{questions} 
        %---------------------------------------------------------------------
	\question Abstract, as written, does not have any actual results. It can be put together without the rest of the paper. Consider including some major results in the abstract, not just final claims. 
	 \begin{solution}
                Thank you for the comment. We have added text to the abstract 
                to discuss some of the results in more concrete detail that 
                support some of the conclusions in the abstract. 
	 \end{solution}


        %---------------------------------------------------------------------
	\question Conclusions consist of fairly detailed discussions. It is recommended
	to include some quantitative metrics to further demonstrate points authors are
	making. The last paragraph of the conclusions needs support by a quantitative
	metric to support authors claim of no significant effect. Overall, each
	conclusions claim needs stronger connection to the results of the paper. 
	 \begin{solution}
	 \end{solution}
 

        %---------------------------------------------------------------------
        \question Pebble Bed Reactor Model (Section II.A). While shortcomings of the
        base model have been briefly discussed, there is no discussion how those
        shortcomings relate to the purposes of the work presented in this paper. The
        overall selected modeling approaches as described in Section II.A appear to be
        overly simplistic and require reactor physics justification of their fidelity
        sufficiency to meet the purpose of this paper and validity of the conclusions.
        Today a wide variety of publications can be found that offer a much more
        detailed representation of the pebble bed reactor physics. Use of the model
        described in Section II.A poses a concern of smearing the results and not
        having enough fidelity to support conclusions and the analysis. Authors are
        commended to review models as the one presented in the paper by Mulyana and
        Chirayath\cite{mulyana_impact_2022}.
        The selected reactor physics metrics need that level of fidelity in
        performance models to capture the details as this study is looking into. If used model is sufficient, that sufficiency must be proven and provided in the paper.
        \begin{solution}
                Thank you for the comment and the suggested reading. We have added
                text to Section II.A to more thoroughly discuss the effects of 
                the pebble homogenization on the results of this study. A more 
                detailed analysis using heterogeneous pebbles would increase 
                the fidelity of the results but was outside the scope of this work. 
                We have expanded the discussion in Section IV.A on how using 
                a heterogeneous pebble model would affect the results and why 
                that future work is of value. 
        \end{solution}
         

        %---------------------------------------------------------------------
        \question The MMR model of Section IIB appears to have the needed fidelity. The
        fidelity of the model with respect to the performance metrics needs to be
        discussed as no such discussion is provided. This reviewer is concerned about
        the noted 2 years burnup step for the 20 years lifetime. This is a rather large
        step that needs justification. With such a large step,  it is expected that the
        depletion model does not capture proper reactor physics of the core during
        operation. The fuel composition and effects cannot be correct with such time
        steps. 
        \begin{solution}
                Thank you for the comment. We updated the 
                burnup step sizes for the MMR depletion 
                to start at 1 day time steps, then increase
                to 5 day steps at 5 days, 10 day steps at 40 days,  
                25 day steps at 100 days,
                110 day steps at 400 days, and then
                to 2 year time steps at 730 days (730 days). 
                These burnup step sizes are based on recommendations from 
                \cite{walter_adaptive_2016}. These step sizes create a finer 
                mesh at the beginning of reactor operation to better 
                capture the core composition changes. 

                We have updated the text in Section II.B to reflect the changes 
                in the burnup step sizes. 
        \end{solution} 


        %---------------------------------------------------------------------
        \question What makes pebble bed and prismatic models and the manner they used
        comparable to be used in the same study? Authors need to offer a discussion of
        what made this paper a single study drawing conclusions for both types of HTREs
        and hence for all HTRs by association. 
        \begin{solution}
        \end{solution} 
         
        %---------------------------------------------------------------------
        \question It would greatly help supporting validity of this work, if authors
        took the time to evaluate reactor physics effects of compositions at the fuel
        lattice level. At that level, pebble bed and prismatic reactors exabit common
        neutronics and results would be very conclusive. 
        \begin{solution}
        \end{solution} 

        %---------------------------------------------------------------------
        \question Authors picked integral and integral differential reactor physics
        characteristics as metrics for the present study. A discussion of these metrics
        and their mathematical nature on the results and their sensitivity to
        composition should be provided. The nature of the selected metrics reduces the
        ability to resolve differences and makes evaluations less sensitive to
        composition effects. This suggests to readers that conclusions of the paper are
        due to metrics and model limitations and not representative of the actual
        reactor physics and fuel effects. If that is not the case, authors need to
        include demonstrations proving the case. The infinite lattice studies and use
        of reaction rates are customary tools for that. 
        \begin{solution}
        \end{solution} 
         
        %---------------------------------------------------------------------
        \question In a study like the one in this paper, it is customary to include a
        sensitivity analysis of the reactor configurations with respect to input
        parameters of interest for the selected metrics. This study is not included but
        is highly recommended.
        \begin{solution}
                Thank you for your comment. This 
                manuscript presents a sensitivity analysis, 
                discussing the sensitivity of reactor 
                physics parameters to the uranium isotopic 
                composition in the fuel. Albeit, it is a 
                coarse sensitivity study given that we only 
                used three different isotopic compositions. 
                We recognize the value in
                sensitivity analysis and therefore have 
                added text to the Conclusions (Section IV)
                to discuss performing a sensitivity 
                analysis with a finer grid space of fuel 
                compositions as recommended future work.

        \end{solution} 


 

\end{questions}
\bibliographystyle{unsrt}
\bibliography{revise}
\end{document}

  %
  % End of file `elsarticle-template-num.tex'.
